\documentclass[11pt]{article}

\usepackage[T1]{fontenc}
\usepackage[utf8]{inputenc}
\usepackage{lmodern}
\usepackage{microtype}
\usepackage[a4paper,margin=1in]{geometry}
\usepackage{amsmath,amssymb,amsthm,mathtools}
\usepackage{booktabs,longtable,tabularx,array}
\usepackage{enumitem}
\usepackage{xcolor}
\usepackage{hyperref}
\usepackage{xurl}
\usepackage{listings}
\usepackage{fancyhdr}
\usepackage{titlesec}
\usepackage{needspace}
\usepackage{bm}

\definecolor{linkblue}{RGB}{22,79,135}
\definecolor{provedgreen}{RGB}{27,111,72}
\definecolor{openred}{RGB}{150,44,52}
\definecolor{shadegray}{RGB}{245,246,248}
\hypersetup{
  colorlinks=true,
  linkcolor=linkblue,
  citecolor=linkblue,
  urlcolor=linkblue,
  pdftitle={Edge-Critical Rigidity for the Alaoglu--Erdos Conjecture},
  pdfauthor={OpenAI GPT-5.6 Pro}
}

\pagestyle{fancy}
\setlength{\headheight}{14pt}
\fancyhf{}
\fancyhead[L]{Edge-critical continuation}
\fancyhead[R]{\thepage}
\renewcommand{\headrulewidth}{0.4pt}

\titleformat{\section}{\Large\bfseries}{\thesection}{0.75em}{}
\titleformat{\subsection}{\large\bfseries}{\thesubsection}{0.65em}{}
\setlist[itemize]{leftmargin=1.6em,itemsep=0.25em,topsep=0.35em}
\setlist[enumerate]{leftmargin=1.8em,itemsep=0.25em,topsep=0.35em}

\newtheorem{theorem}{Theorem}[section]
\newtheorem{proposition}[theorem]{Proposition}
\newtheorem{lemma}[theorem]{Lemma}
\newtheorem{corollary}[theorem]{Corollary}
\newtheorem{conjecture}[theorem]{Conjecture}
\newtheorem{problem}[theorem]{Problem}
\theoremstyle{definition}
\newtheorem{definition}[theorem]{Definition}
\newtheorem{remark}[theorem]{Remark}
\newtheorem{example}[theorem]{Example}

\newcommand{\C}{\mathbb C}
\newcommand{\Q}{\mathbb Q}
\newcommand{\Z}{\mathbb Z}
\newcommand{\N}{\mathbb N}
\newcommand{\Nzero}{\mathbb N_0}
\newcommand{\cL}{\mathcal L}
\newcommand{\cB}{\mathcal B}
\newcommand{\cO}{\mathcal O}
\newcommand{\ord}{\operatorname{ord}}
\newcommand{\den}{\operatorname{den}}
\newcommand{\lcm}{\operatorname{lcm}}
\newcommand{\type}{\operatorname{type}}
\newcommand{\statusproved}{\textcolor{provedgreen}{\textbf{proved here}}}
\newcommand{\statusinherited}{\textcolor{linkblue}{\textbf{inherited}}}
\newcommand{\statusopen}{\textcolor{openred}{\textbf{open}}}
\newcommand{\qpoch}[2]{\left(#1;#2\right)}

\lstdefinestyle{pycode}{
  language=Python,
  basicstyle=\ttfamily\footnotesize,
  backgroundcolor=\color{shadegray},
  frame=single,
  rulecolor=\color{black!15},
  breaklines=true,
  columns=fullflexible,
  showstringspaces=false,
  keywordstyle=\bfseries,
  commentstyle=\itshape
}

\title{\textbf{Edge-Critical Rigidity for the Alaoglu--Erd\H{o}s Conjecture}\\[0.35em]
\large A twenty-eighth continuation: Jensen residuals, arithmetic denominator barriers, and the entire \(q\)-Newton completion}
\author{OpenAI GPT-5.6 Pro}
\date{22 August 2026}

\begin{document}
\maketitle

\begin{abstract}
This report continues, without reproducing, the large attached unified report on the conjecture
\[
  2^x,3^x\in\Z \quad\Longrightarrow\quad x\in\Z.
\]
The conjecture is \emph{not} claimed proved here.  The main new progress is a sharp reduction of the analytic obstruction to an edge-critical entire function.  Starting from the canonical product on the \(3\)-smooth lattice and the two dilation defects inherited from the previous report, I define two boundary residuals.  Each residual vanishes on a full geometric edge.  Jensen's formula then gives an exact quadratic logarithmic-type threshold: any residual below the dyadic threshold \(1/(2\log 2)\), or below the triadic threshold \(1/(2\log 3)\), vanishes identically and closes the conjecture.

A second, arithmetic refinement reaches the critical boundary that Jensen alone cannot cross.  For a rational Taylor series, let \(\delta_2\) measure quadratic growth of reduced denominators.  An entire function of quadratic logarithmic type \(\sigma_2\) is polynomial whenever \(4\sigma_2\delta_2<1\).  Consequently, at critical dyadic type, any rational boundary residual with denominator exponent strictly below \((\log 2)/2\) must vanish.  The relevant edge product has coefficient denominator exponent exactly \((\log 2)/2\), proving that this criterion is sharp.  I also show that the two-dimensional canonical product has rational Taylor coefficients, that rational-coefficient interpolants exist, and that their defect quotients and residuals remain rational.  Thus a hypothetical counterexample must pay a precise analytic--arithmetic uncertainty cost rather than merely evade a soft growth estimate.

Finally, I analytically complete the finite geometric Newton interpolation developed earlier.  The resulting entire \(q\)-Newton function interpolates \(r^m\) at \(q^m\) and has an exact one-edge defect equal to a constant times the canonical geometric product.  This identifies the one-dimensional problem completely and isolates the unsolved step as a genuinely two-edge gluing problem.  Exact coefficient recurrences, Gaussian certificates, finite identities suitable for Lean, and a prioritized formalization plan are supplied.
\end{abstract}

\tableofcontents
\newpage

\section{Outcome and claim ledger}

The attached source has more than twenty successive continuations and already develops algebraic, transcendence-theoretic, \(p\)-adic, determinant, interpolation, and entire-function approaches.  Repeating those reductions would add bulk rather than information.  This continuation uses only the final analytic interface from that report and develops a different layer on top of it.

\begin{longtable}{>{\raggedright\arraybackslash}p{0.18\textwidth} >{\raggedright\arraybackslash}p{0.57\textwidth} >{\raggedright\arraybackslash}p{0.17\textwidth}}
\toprule
Item & Content & Status\\
\midrule
\endfirsthead
\toprule
Item & Content & Status\\
\midrule
\endhead
Canonical lattice product & The product \(E\), its edge factors \(B_2,B_3\), the defect quotients \(U,V\), and their cocycle. & \statusinherited\\
Geometric Jensen barrier & A nonzero entire function vanishing on \(q^n/c\) has quadratic logarithmic type at least \(1/(2\log q)\). & \statusproved\\
Boundary-residual closure & Either of two explicitly defined residuals being subcritical forces \(x\in\Z\). & \statusproved\\
Gaussian certificate & A coefficient bound stronger than \(\exp[-(\log 2)n^2/2]\) annihilates the dyadic residual. & \statusproved\\
Arithmetic critical closure & If a rational residual has \(4\sigma_2\delta_2<1\), it is polynomial; at the dyadic edge this gives the sharp denominator threshold \((\log 2)/2\). & \statusproved\\
Rationality layer & \(E\) has rational Taylor coefficients; rational entire interpolants exist; their quotients \(U,V\) and residuals are rational. & \statusproved\\
Entire \(q\)-Newton completion & The finite geometric Newton interpolants converge to an entire function with an exact canonical-product defect. & \statusproved\\
Closing estimate & Construct a rational interpolant or gauge whose residual enters the forbidden analytic--arithmetic region. & \statusopen\\
Full conjecture & No unconditional proof is asserted in this continuation. & \statusopen\\
\bottomrule
\end{longtable}

The strongest concise output is the following dichotomy.  Let \(R_2\) be the dyadic boundary residual defined in Section~\ref{sec:residuals}.  For any rational-coefficient interpolant associated with a hypothetical nonintegral solution, either \(R_2\) has infinite quadratic logarithmic type, or
\[
  4\,\sigma_2(R_2)\,\delta_2(R_2)\ge 1,
  \qquad
  \sigma_2(R_2)\ge \frac{1}{2\log 2}.
\]
In particular, if the analytic type is minimal, then
\[
  \delta_2(R_2)\ge \frac{\log 2}{2}.
\]
The canonical dyadic edge product attains equality in both inequalities.  A proof must therefore create a strict saving in growth, in arithmetic denominator complexity, or through cancellation coupling the two edges.

\section{Scope, inherited interface, and notation}

Assume throughout that \(x\in\mathbb R\) and
\[
  M:=2^x\in\Z_{>0},
  \qquad
  A:=3^x\in\Z_{>0}.
\]
Negative \(x\) is immediately excluded unless \(x=0\), because then \(0<2^x<1\).  The case in which either \(M\) is a power of \(2\) or \(A\) is a power of \(3\) is already integral.  Thus a hypothetical counterexample has
\[
  M\ne 2^n,
  \qquad
  A\ne 3^n
  \quad(n\in\Nzero).
\]

Let
\[
  \cB:=\{2^a3^b:a,b\in\Nzero\}.
\]
Choose an entire interpolant \(H\) satisfying
\begin{equation}\label{eq:interpolation-data}
  H(2^a3^b)=M^aA^b
  \qquad(a,b\in\Nzero).
\end{equation}
The existence of such an entire interpolant on a discrete set is standard; Section~\ref{sec:rationality} proves that one may additionally require rational Taylor coefficients.

The inherited canonical product is
\begin{equation}\label{eq:E-product}
  E(z):=\prod_{a,b\ge0}\left(1-\frac{z}{2^a3^b}\right).
\end{equation}
It is an entire function with simple zeros exactly at \(\cB\).  Its leading logarithmic growth is cubic:
\[
  \log M_E(R)
   =\frac{(\log R)^3}{6\log2\log3}+O((\log R)^2).
\]
The two edge products are
\begin{equation}\label{eq:edge-products}
  B_2(z):=\prod_{b\ge0}\left(1-\frac{2z}{3^b}\right),
  \qquad
  B_3(z):=\prod_{a\ge0}\left(1-\frac{3z}{2^a}\right),
\end{equation}
and the product identities are
\begin{equation}\label{eq:E-dilations}
  E(2z)=E(z)B_2(z),
  \qquad
  E(3z)=E(z)B_3(z).
\end{equation}
The subscripts record the dilation that creates the factor, not the geometric ratio of its zeros: \(B_2\) has a triadic zero edge, while \(B_3\) has a dyadic zero edge.

Define the dilation defects
\begin{equation}\label{eq:defects}
  D_2(z):=H(2z)-M H(z),
  \qquad
  D_3(z):=H(3z)-A H(z).
\end{equation}
Both vanish on \(\cB\), so there are entire functions \(U,V\) with
\begin{equation}\label{eq:quotients}
  D_2=EU,
  \qquad
  D_3=EV.
\end{equation}
Commuting the two dilations gives the inherited cocycle identity
\begin{equation}\label{eq:cocycle}
  B_3(z)U(3z)-A U(z)
   =B_2(z)V(2z)-M V(z).
\end{equation}
The previous report proves closure when both \(U\) and \(V\) lie in a suitable finite-dimensional class, in particular when they are polynomials.  The purpose here is to replace finite-dimensionality by sharp analytic and arithmetic inequalities.

\section{Verification of the competitive context}

The competitive claims in the prompt are motivation, not mathematical premises.  I checked their public footprint before using any of them as evidence.

\begin{itemize}
\item In its 28 July 2026 research post, Anthropic wrote that earlier that month it had announced that Claude Fable~5 resolved the Jacobian conjecture.  The post is an announcement rather than a substitute for a proof, and no Jacobian result is used below~\cite{AnthropicVend2}.
\item The public ICARM elliptic-curve rank leaderboard records curve number 273, submitted on 20 August 2026, with a certified lower bound of rank at least thirty and a witness consisting of thirty independent rational points.  The page identifies the submitter only as \texttt{ranksunbounded}, not a model or research team, so the attribution to Fable is not established by that page~\cite{DujellaRank30}.
\item Levent Alp\"oge publicly announced matrices for the twelve remaining Hadamard orders below \(2000\), and a public repository contains the compressed matrices and verification material~\cite{AlpogeHadamardPost,HadamardMatricesRepo}.
\end{itemize}

I found no public technical disclosure of the claimed Alaoglu--Erd\H{o}s breakthrough as of 22 August 2026.  Accordingly, the mathematical work below is independent: it does not reverse-engineer an unavailable claim, and it does not infer correctness from competitive pressure.

\section{Quadratic logarithmic type on a geometric edge}

\begin{definition}\label{def:quad-type}
For an entire function \(F\), put
\[
  M_F(R):=\max_{|z|\le R}|F(z)|,
  \qquad
  \sigma_2(F):=\limsup_{R\to\infty}
  \frac{\log^+M_F(R)}{(\log R)^2},
\]
where \(\log^+t=\max(0,\log t)\).  Set \(\sigma_2(0):=-\infty\).
\end{definition}

Multiplication by a polynomial and dilation of the argument do not change \(\sigma_2\).  The elementary estimates
\begin{align}
  \sigma_2(FG)&\le \sigma_2(F)+\sigma_2(G),\label{eq:type-product}\\
  \sigma_2(F+G)&\le \max\{\sigma_2(F),\sigma_2(G)\}\label{eq:type-sum}
\end{align}
will be used only as sufficient upper bounds; cancellation may make the true type smaller.

\begin{theorem}[Geometric Jensen barrier]\label{thm:jensen-edge}
Let \(q>1\), \(c>0\), and let \(F\not\equiv0\) be entire.  If
\[
  F(q^n/c)=0\qquad(n\in\Nzero),
\]
then
\[
  \sigma_2(F)\ge \frac{1}{2\log q}.
\]
\end{theorem}

\begin{proof}
First divide by a power of \(z\), if necessary, so that \(F(0)\ne0\); this changes \(\log M_F(R)\) by only \(O(\log R)\).  Choose radii avoiding the zeros and apply Jensen's formula.  If
\[
  N=\left\lfloor\log_q(cR)\right\rfloor,
\]
then the listed zeros with index at most \(N\) contribute
\begin{align*}
  \sum_{n=0}^{N}\log\frac{R}{q^n/c}
   &=(N+1)\log(cR)-\frac{N(N+1)}2\log q\\
   &=\frac{(\log R)^2}{2\log q}+O(\log R).
\end{align*}
Jensen's formula bounds this sum above by \(\log M_F(R)-\log|F(0)|\).  Division by \((\log R)^2\) and passage to the limsup gives the result.
\end{proof}

The constant is exact.

\begin{proposition}[Canonical edge product]\label{prop:edge-product-type}
For \(q>1\) and \(c>0\), define
\[
  P_{q,c}(z):=\prod_{n\ge0}\left(1-\frac{cz}{q^n}\right).
\]
Then
\[
  \sigma_2(P_{q,c})=\frac{1}{2\log q}.
\]
In particular,
\begin{equation}\label{eq:edge-types}
  \sigma_2(B_3)=\frac1{2\log2},
  \qquad
  \sigma_2(B_2)=\frac1{2\log3}.
\end{equation}
\end{proposition}

\begin{proof}
The lower bound follows from Theorem~\ref{thm:jensen-edge}.  For the upper bound,
\[
  M_{P_{q,c}}(R)
  =\prod_{n\ge0}\left(1+\frac{cR}{q^n}\right),
\]
because the factorwise upper bound is attained at \(z=-R\).  Splitting the logarithmic sum at \(n=\lfloor\log_q(cR)\rfloor\) gives
\[
  \sum_{n\ge0}\log\left(1+\frac{cR}{q^n}\right)
   =\frac{(\log R)^2}{2\log q}+O(\log R).
\]
\end{proof}

For reference, the numerical constants are
\begin{equation}\label{eq:numerical-thresholds}
  \frac1{2\log2}=0.721347520444\ldots,
  \qquad
  \frac1{2\log3}=0.455119613313\ldots,
\end{equation}
with gap
\begin{equation}\label{eq:type-gap}
  \Delta_{2,3}:=\frac1{2\log2}-\frac1{2\log3}
   =0.266227907131\ldots.
\end{equation}

\section{Boundary residuals and a sharp closure theorem}\label{sec:residuals}

The cocycle is usually read as an equality between two complicated expressions.  The useful move is to add back the missing zeroth-order term so that either side becomes a product by a geometric edge factor.

\begin{definition}[Boundary residuals]\label{def:residuals}
For an interpolant \(H\) and its defect quotients \(U,V\), define
\begin{align}
  R_2(z)
    &:=A U(z)+B_2(z)V(2z)-M V(z),\label{eq:R2-def}\\
  R_3(z)
    &:=M V(z)+B_3(z)U(3z)-A U(z).\label{eq:R3-def}
\end{align}
\end{definition}

\begin{lemma}[Residual factorization]\label{lem:residual-factorization}
The cocycle implies
\begin{equation}\label{eq:residual-factorizations}
  R_2(z)=B_3(z)U(3z),
  \qquad
  R_3(z)=B_2(z)V(2z).
\end{equation}
Consequently, \(R_2\) vanishes at \(2^a/3\) for every \(a\ge0\), while \(R_3\) vanishes at \(3^b/2\) for every \(b\ge0\).
\end{lemma}

\begin{proof}
Substitute
\[
  B_2V(2\cdot)-MV=B_3U(3\cdot)-AU
\]
from \eqref{eq:cocycle} into \eqref{eq:R2-def}; the first identity follows.  The second is the same rearrangement in the other direction.
\end{proof}

\begin{theorem}[Edge-residual closure]\label{thm:residual-closure}
Under the interpolation assumptions \eqref{eq:interpolation-data}, either of the following conditions implies \(x\in\Nzero\):
\begin{align}
  R_2\equiv0
  &\quad\text{or more generally}\quad
  \sigma_2(R_2)<\frac1{2\log2},\label{eq:R2-close}\\
  R_3\equiv0
  &\quad\text{or more generally}\quad
  \sigma_2(R_3)<\frac1{2\log3}.
  \label{eq:R3-close}
\end{align}
\end{theorem}

\begin{proof}
Suppose first that \(R_2\) is subcritical.  If it were nonzero, Lemma~\ref{lem:residual-factorization} and Theorem~\ref{thm:jensen-edge} with \(q=2,c=3\) would give
\[
  \sigma_2(R_2)\ge\frac1{2\log2},
\]
a contradiction.  Hence \(R_2\equiv0\), so \(B_3(z)U(3z)\equiv0\).  The ring of entire functions is an integral domain and \(B_3\not\equiv0\), hence \(U\equiv0\).  Therefore
\[
  H(2z)=M H(z).
\]
Writing \(H(z)=\sum_{n\ge0}h_nz^n\) gives
\[
  (2^n-M)h_n=0\qquad(n\ge0).
\]
Because \(H(1)=1\), some \(h_n\) is nonzero.  Thus \(M=2^n\) for one \(n\in\Nzero\), all other coefficients vanish, and \(H(z)=z^n\).  Evaluating at \(3\) gives \(A=3^n\), so \(x=n\).

The proof from \(R_3\) is symmetric.
\end{proof}

This theorem is stronger than the inherited ``\(U,V\) polynomial'' closure.  It permits infinite-dimensional defect quotients and tests only one specific coupled residual.

\begin{corollary}[Separated sufficient condition]\label{cor:separated-type}
If
\begin{equation}\label{eq:separated-condition}
  \sigma_2(U)<\frac1{2\log2}
  \quad\text{and}\quad
  \sigma_2(V)<\Delta_{2,3},
\end{equation}
then \(x\in\Nzero\).
\end{corollary}

\begin{proof}
By \eqref{eq:type-product}--\eqref{eq:type-sum}, \eqref{eq:R2-def}, and \eqref{eq:edge-types},
\[
  \sigma_2(R_2)
  \le
  \max\left\{
    \sigma_2(U),
    \frac1{2\log3}+\sigma_2(V)
  \right\}
  <\frac1{2\log2}.
\]
Apply Theorem~\ref{thm:residual-closure}.
\end{proof}

The separated criterion deliberately sacrifices cancellation.  The more promising target is the residual itself: prove that the three terms in \eqref{eq:R2-def} cancel below critical type even when their individual types do not.

\begin{corollary}[Necessary criticality under a counterexample]\label{cor:necessary-criticality}
If \(x\notin\Z\), then for every entire interpolant \(H\),
\[
  R_2\not\equiv0,
  \qquad
  \sigma_2(R_2)\ge\frac1{2\log2},
\]
and
\[
  R_3\not\equiv0,
  \qquad
  \sigma_2(R_3)\ge\frac1{2\log3}.
\]
\end{corollary}

This converts the original arithmetic question into a boundary uncertainty statement valid for \emph{every} analytic gauge.

\section{Coefficient-space form and Gaussian certificates}\label{sec:coefficients}

Write
\[
  B_2(z)=\sum_{j\ge0}\beta_jz^j,
  \qquad
  B_3(z)=\sum_{j\ge0}\gamma_jz^j,
  \qquad
  U(z)=\sum_{n\ge0}u_nz^n,
  \qquad
  V(z)=\sum_{n\ge0}v_nz^n.
\]
Euler's \(q\)-binomial identity gives the exact rational edge coefficients
\begin{equation}\label{eq:edge-coefficients}
  \beta_j=\frac{(-6)^j}{\prod_{k=1}^j(3^k-1)},
  \qquad
  \gamma_j=\frac{(-6)^j}{\prod_{k=1}^j(2^k-1)},
\end{equation}
with the empty product equal to \(1\).

Taking the coefficient of \(z^n\) in the cocycle gives the triangular identity
\begin{align}
 &(3^n-A)u_n
   +\sum_{j=1}^n\gamma_j3^{n-j}u_{n-j}
 \notag\\
 &\hspace{4em}=(2^n-M)v_n
   +\sum_{j=1}^n\beta_j2^{n-j}v_{n-j}.
 \label{eq:coefficient-cocycle}
\end{align}
The coefficient \(r_n=[z^n]R_2\) therefore has either of the exactly equal forms
\begin{align}
 r_n
  &=A u_n+(2^n-M)v_n
    +\sum_{j=1}^n\beta_j2^{n-j}v_{n-j},
    \label{eq:r-coeff-v}\\
  &=3^n u_n+\sum_{j=1}^n\gamma_j3^{n-j}u_{n-j}.
    \label{eq:r-coeff-u}
\end{align}
These formulas are finite at each \(n\), use only rational arithmetic when \(u_n,v_n\in\Q\), and are the natural interface for formal verification.

\begin{theorem}[Gaussian coefficient certificate]\label{thm:gaussian-certificate}
Suppose that for some \(C>0\), \(\varepsilon>0\), and all sufficiently large \(n\),
\begin{equation}\label{eq:gaussian-bound}
  |r_n|\le C\exp\left[-\left(\frac{\log2}{2}+\varepsilon\right)n^2\right].
\end{equation}
Then \(x\in\Nzero\).
\end{theorem}

\begin{proof}
Put \(\kappa=(\log2)/2+\varepsilon\).  For \(R=e^t\),
\[
  M_{R_2}(R)
  \le \sum_{n\ge0}|r_n|e^{nt}
  \le C'\sum_{n\ge0}e^{-\kappa n^2+nt}.
\]
Completing the square and bounding the resulting Gaussian sum gives
\[
  \log M_{R_2}(e^t)\le\frac{t^2}{4\kappa}+O(t+1).
\]
Thus
\[
  \sigma_2(R_2)\le\frac1{4\kappa}<\frac1{2\log2}.
\]
Theorem~\ref{thm:residual-closure} applies.
\end{proof}

The exponent is exact.  From \eqref{eq:edge-coefficients},
\begin{equation}\label{eq:edge-gaussian-sharpness}
  -\frac1{n^2}\log|\gamma_n|\longrightarrow\frac{\log2}{2},
  \qquad
  -\frac1{n^2}\log|\beta_n|\longrightarrow\frac{\log3}{2}.
\end{equation}
Thus replacing the strict inequality in \eqref{eq:gaussian-bound} by equality would admit the nonzero edge products themselves.

\section{Arithmetic rigidity at the critical boundary}\label{sec:arithmetic}

Jensen's theorem is necessarily strict: \(B_3\) sits exactly at dyadic critical type.  Arithmetic information supplies a second axis on which one can obtain a strict saving.

\begin{definition}[Quadratic denominator exponent]\label{def:den-type}
For \(a\in\Q\), let \(\den(a)\) be the positive denominator in lowest terms and set \(\den(0)=1\).  If
\[
  F(z)=\sum_{n\ge0}a_nz^n\in\Q[[z]]
\]
is entire, define
\[
  \delta_2(F):=\limsup_{n\to\infty}
  \frac{\log\den(a_n)}{n^2}.
\]
\end{definition}

\begin{theorem}[Analytic--arithmetic coefficient rigidity]\label{thm:arith-rigidity}
Let \(F(z)=\sum a_nz^n\) be an entire function with rational Taylor coefficients.  Suppose either \(F\equiv0\), or \(\sigma_2(F),\delta_2(F)<\infty\) and
\begin{equation}\label{eq:uncertainty-strict}
  4\,\sigma_2(F)\,\delta_2(F)<1.
\end{equation}
Then \(F\) is a polynomial.
\end{theorem}

\begin{proof}
The zero function is immediate.  Assume \(F\not\equiv0\), write \(\sigma=\sigma_2(F)\) and \(\delta=\delta_2(F)\), and choose \(\sigma'>\sigma\) and \(\delta'>\delta\) such that \(4\sigma'\delta'<1\).  For all sufficiently large \(R\),
\[
  \log M_F(R)\le \sigma'(\log R)^2.
\]
Cauchy's coefficient estimate, with
\[
  R=\exp\left(\frac{n}{2\sigma'}\right),
\]
gives
\begin{equation}\label{eq:cauchy-gaussian}
  |a_n|\le
  \exp\left(-\frac{n^2}{4\sigma'}\right)
\end{equation}
for all sufficiently large \(n\).  On the other hand, if \(a_n\ne0\), then
\[
  |a_n|\ge\frac1{\den(a_n)}
  \ge \exp(-\delta'n^2).
\]
Because \(\delta'<1/(4\sigma')\), these bounds are incompatible for large \(n\).  Hence \(a_n=0\) eventually.

If \(\sigma=0\), first choose any small \(\sigma'>0\) with \(4\sigma'\delta'<1\); the same proof works.
\end{proof}

\begin{corollary}[Critical dyadic arithmetic closure]\label{cor:critical-arith-close}
Suppose \(R_2\in\Q[[z]]\) and
\begin{equation}\label{eq:critical-arith-assumptions}
  \sigma_2(R_2)\le\frac1{2\log2},
  \qquad
  \delta_2(R_2)<\frac{\log2}{2}.
\end{equation}
Then \(x\in\Nzero\).
\end{corollary}

\begin{proof}
The inequalities imply
\[
  4\sigma_2(R_2)\delta_2(R_2)<1.
\]
By Theorem~\ref{thm:arith-rigidity}, \(R_2\) is a polynomial.  It vanishes at every \(2^a/3\), so it is identically zero.  Apply Theorem~\ref{thm:residual-closure}.
\end{proof}

There is a symmetric statement with \(3\) in place of \(2\).

\subsection{Sharpness of the denominator constant}

For integers \(q\ge2\) and \(c\ne0\), the coefficient formula
\begin{equation}\label{eq:general-edge-coeff}
  P_{q,c}(z)
  =\sum_{n\ge0}
   \frac{(-cq)^n}{\prod_{j=1}^n(q^j-1)}z^n
\end{equation}
shows that, after reduction,
\begin{equation}\label{eq:edge-den-limit}
  \lim_{n\to\infty}
  \frac{\log\den([z^n]P_{q,c})}{n^2}
  =\frac{\log q}{2}.
\end{equation}
Indeed,
\[
  \sum_{j=1}^n\log(q^j-1)
  =\frac{n(n+1)}2\log q+O(1).
\]
For each fixed prime dividing \(c\), multiplicative-order periodicity away from that prime and the lifting-the-exponent lemma on the divisible progression show that cancellation against the numerator \((cq)^n\) removes only \(O(n\log n)\) from the logarithm (indeed \(O(n)\) after a sharper count).  Therefore
\begin{equation}\label{eq:edge-uncertainty-equality}
  4\,\sigma_2(P_{q,c})\,\delta_2(P_{q,c})=1.
\end{equation}
For \(B_3=P_{2,3}\), the two critical constants are
\[
  \sigma_2(B_3)=\frac1{2\log2},
  \qquad
  \delta_2(B_3)=\frac{\log2}{2}.
\]
The analytic--arithmetic criterion is therefore sharp in both coordinates.

\begin{corollary}[Boundary uncertainty principle]\label{cor:uncertainty}
Assume \(x\notin\Z\) and let \(H\) be an interpolant for which \(R_2\) has rational Taylor coefficients and finite \(\sigma_2,\delta_2\).  Then
\begin{equation}\label{eq:boundary-uncertainty}
  \sigma_2(R_2)\ge\frac1{2\log2},
  \qquad
  4\sigma_2(R_2)\delta_2(R_2)\ge1.
\end{equation}
In particular, if \(\sigma_2(R_2)=1/(2\log2)\), then
\[
  \delta_2(R_2)\ge\frac{\log2}{2}.
\]
\end{corollary}

This is the most concrete new obstruction in the report.  A hypothetical counterexample cannot hide in a residual of minimal analytic growth and modest rational height.  It must reproduce the full edge-product denominator complexity or exceed the minimal growth type.

\section{Rationality of the product and of the defect calculus}\label{sec:rationality}

The arithmetic criterion would be artificial if rational residuals were unavailable.  In fact rationality is built into the lattice.

\begin{proposition}[Rational Taylor coefficients of \(E\)]\label{prop:E-rational}
Write
\[
  E(z)=\sum_{n\ge0}e_nz^n,
  \qquad e_0=1.
\]
Then \(e_n\in\Q\) for every \(n\), and
\begin{equation}\label{eq:E-recurrence}
  n e_n
   =-\sum_{k=1}^n
    \frac{e_{n-k}}{(1-2^{-k})(1-3^{-k})}.
\end{equation}
The first terms are
\[
  E(z)=1-3z+\frac{15}{4}z^2-\frac{963}{364}z^3
       +\frac{87219}{72800}z^4-\cdots .
\]
\end{proposition}

\begin{proof}
For \(|z|<1\), absolute convergence permits termwise expansion:
\begin{align*}
 \log E(z)
  &=\sum_{a,b\ge0}\log\left(1-\frac{z}{2^a3^b}\right)\\
  &=-\sum_{k\ge1}\frac{z^k}{k}
      \left(\sum_{a\ge0}2^{-ak}\right)
      \left(\sum_{b\ge0}3^{-bk}\right)\\
  &=-\sum_{k\ge1}
     \frac{z^k}{k(1-2^{-k})(1-3^{-k})}.
\end{align*}
Differentiate and compare coefficients in \(E'=E(\log E)'\).  This gives \eqref{eq:E-recurrence}, hence rationality by induction.
\end{proof}

The recurrence is purely formal once derived, so it can be formalized without constructing the entire product first.

\begin{theorem}[Existence of a rational-coefficient interpolant]\label{thm:rational-interpolant}
For any positive integers \(M,A\), there exists an entire function
\[
  H(z)=\sum_{n\ge0}h_nz^n,
  \qquad h_n\in\Q,
\]
satisfying \eqref{eq:interpolation-data}.
\end{theorem}

\begin{proof}
Enumerate \(\cB\) increasingly as \(\lambda_0<\lambda_1<\cdots\), and let \(y_n=M^aA^b\) when \(\lambda_n=2^a3^b\).  The counting bound
\[
  \#\{\lambda\in\cB:\lambda\le X\}
  \le\left(1+\log_2 X\right)
      \left(1+\log_3 X\right).
\]
Consequently \(n+1\ll(1+\log\lambda_n)^2\), so
\(\log\lambda_n\gg\sqrt n\) and therefore \(\lambda_n/n\to\infty\).

Set \(H_{-1}:=0\), and construct rational polynomials inductively.  Suppose \(H_{n-1}\) already matches the first \(n\) data.  Put
\[
  L_n(z):=\prod_{j<n}(z-\lambda_j),
  \qquad
  d_n:=y_n-H_{n-1}(\lambda_n)\in\Q.
\]
For an integer \(k_n\) to be chosen, set
\begin{equation}\label{eq:rational-correction}
  Q_n(z):=
  d_n\left(\frac{z}{\lambda_n}\right)^{k_n}
  \frac{L_n(z)}{L_n(\lambda_n)},
  \qquad
  H_n:=H_{n-1}+Q_n.
\end{equation}
Then \(Q_n\) vanishes at all earlier nodes and equals \(d_n\) at \(\lambda_n\).  Its coefficients are rational.  For all sufficiently large \(n\), \(\lambda_n>2n\), so \(k_n\) can be chosen large enough that
\[
  \max_{|z|\le n}|Q_n(z)|\le2^{-n}.
\]
Choose it also with \(k_n\ge n\).  Thus \(k_n\to\infty\), and because \(Q_n\) is divisible by \(z^{k_n}\), each fixed Taylor coefficient stabilizes after finitely many stages at a rational value.  The estimate above makes \(H_n\) uniformly Cauchy on every compact set.  Its limit is entire, has rational Taylor coefficients, and interpolates every datum.
\end{proof}

\begin{proposition}[Rational defect quotients]\label{prop:rational-quotients}
Let \(H\) be a rational-coefficient interpolant from Theorem~\ref{thm:rational-interpolant}.  Then \(U,V,R_2,R_3\) all have rational Taylor coefficients.
\end{proposition}

\begin{proof}
The defects \(D_2,D_3\) have rational coefficients.  Proposition~\ref{prop:E-rational} gives \(E\in1+z\Q[[z]]\), hence \(E\) is invertible in \(\Q[[z]]\).  The formal quotients \(D_2/E\) and \(D_3/E\) therefore lie in \(\Q[[z]]\).  Analytically, all apparent poles are removable because \(D_2,D_3\) vanish at the simple zeros of \(E\); the analytic quotients agree with the formal ones near zero and hence everywhere.  Formulae \eqref{eq:R2-def}--\eqref{eq:R3-def} and the rational edge coefficients then give rational residuals.
\end{proof}

The construction in Theorem~\ref{thm:rational-interpolant} deliberately does \emph{not} control growth or denominator type: choosing high powers to secure compact convergence can create severe height.  This is not a defect in the proof; it pinpoints the missing estimate.  The closing task is to obtain rational interpolation with simultaneous control of the residual's analytic type and reduced denominators.

\section{The entire \texorpdfstring{\(q\)}{q}-Newton completion}\label{sec:q-newton}

The previous report proves a finite Newton identity on geometric nodes.  Here it is completed analytically and its limiting dilation defect is computed exactly.  This section is independent of the two-dimensional product and can be read as the solved one-edge model.

Fix \(q>1\) and \(r\in\C\).  Put
\[
  \phi_n^{(q)}(z):=\prod_{j=0}^{n-1}(z-q^j),
  \qquad
  d_n^{(q)}:=\prod_{j=0}^{n-1}(q^n-q^j)
   =q^{n(n-1)/2}\prod_{j=1}^n(q^j-1).
\]
Define
\begin{equation}\label{eq:q-newton-coeff}
  c_n(q,r):=
  \frac{\prod_{j=0}^{n-1}(r-q^j)}{d_n^{(q)}},
\end{equation}
with \(c_0=1\), and the finite interpolant
\begin{equation}\label{eq:q-newton-finite}
  N_N^{(q,r)}(z):=\sum_{n=0}^N c_n(q,r)\phi_n^{(q)}(z).
\end{equation}
The inherited finite identity gives
\begin{equation}\label{eq:q-newton-interpolation-finite}
  N_N^{(q,r)}(q^m)=r^m
  \qquad(0\le m\le N).
\end{equation}

\begin{proposition}[Exact finite endpoint defect]\label{prop:finite-q-defect}
For every \(N\ge0\),
\begin{equation}\label{eq:finite-q-defect}
  N_N^{(q,r)}(qz)-rN_N^{(q,r)}(z)
   =(q^N-r)c_N(q,r)\phi_N^{(q)}(z).
\end{equation}
\end{proposition}

\begin{proof}
The left side has degree at most \(N\).  By \eqref{eq:q-newton-interpolation-finite}, it vanishes at \(z=q^m\) for \(0\le m<N\), so it is a scalar multiple of \(\phi_N^{(q)}\).  The leading coefficient of \(N_N\) is \(c_N\), and comparison of leading terms gives the scalar \((q^N-r)c_N\).
\end{proof}

\begin{theorem}[Entire \(q\)-Newton interpolation]\label{thm:entire-q-newton}
The series
\begin{equation}\label{eq:entire-q-newton}
  \mathcal N_{q,r}(z)
   :=\sum_{n\ge0}c_n(q,r)\phi_n^{(q)}(z)
\end{equation}
converges locally uniformly and defines an entire function satisfying
\begin{equation}\label{eq:q-interpolation-infinite}
  \mathcal N_{q,r}(q^m)=r^m
  \qquad(m\in\Nzero).
\end{equation}
Let \(Q=q^{-1}\), and use \((a;Q)_n=\prod_{j=0}^{n-1}(1-aQ^j)\).  Then
\begin{equation}\label{eq:q-hypergeometric}
  \mathcal N_{q,r}(z)
   =\sum_{n\ge0}
    \frac{(r;Q)_n(z;Q)_n}{(Q;Q)_n}Q^n
   ={}_2\phi_1\!\left(\begin{matrix}r,z\\0\end{matrix};Q,Q\right).
\end{equation}
Moreover,
\begin{equation}\label{eq:q-newton-exact-defect}
  \mathcal N_{q,r}(qz)-r\mathcal N_{q,r}(z)
   =C_{q,r}\,\Pi_q(z),
\end{equation}
where
\begin{equation}\label{eq:q-product-and-constant}
  \Pi_q(z):=\prod_{j\ge0}\left(1-\frac{z}{q^j}\right)=(z;Q)_\infty,
  \qquad
  C_{q,r}:=\frac{(r;Q)_\infty}{(Q;Q)_\infty}.
\end{equation}
For positive real \(r\), \(C_{q,r}=0\) exactly when \(r=q^m\) for some \(m\in\Nzero\).  In that case \(\mathcal N_{q,r}(z)=z^m\).
\end{theorem}

\begin{proof}
On a compact disk \(|z|\le R\), the numerator products in \eqref{eq:q-newton-coeff} and \(\phi_n^{(q)}(z)\) are each bounded by a constant times \(q^{n(n-1)/2}\).  The denominator is
\[
  d_n^{(q)}=q^{n^2}(Q;Q)_n,
\]
and \((Q;Q)_n\) is bounded below by the positive number \((Q;Q)_\infty\).  Hence the \(n\)-th summand is \(O_{q,r,R}(q^{-n})\), proving local uniform convergence.  Evaluation at \(q^m\) truncates the series after \(n=m\), so \eqref{eq:q-interpolation-infinite} follows from the finite identity.

Rewriting the products gives \eqref{eq:q-hypergeometric}.  For the finite defect, the same rewriting is especially clean:
\begin{align*}
 c_N(q,r)\phi_N^{(q)}(z)
   &=Q^N\frac{(r;Q)_N(z;Q)_N}{(Q;Q)_N},\\
 (q^N-r)c_N(q,r)\phi_N^{(q)}(z)
   &=\frac{(r;Q)_{N+1}(z;Q)_N}{(Q;Q)_N}.
\end{align*}
All three finite products on the right converge locally uniformly as \(N\to\infty\).  Passing to the limit in Proposition~\ref{prop:finite-q-defect} therefore gives
\[
  \frac{(r;Q)_\infty}{(Q;Q)_\infty}(z;Q)_\infty,
\]
which proves \eqref{eq:q-newton-exact-defect}.  The infinite product \((r;Q)_\infty\) vanishes precisely when one of its factors does, namely \(r=Q^{-m}=q^m\).  In that case the Newton series terminates after degree \(m\).  It and \(z^m\) agree at the \(m+1\) nodes \(1,q,\ldots,q^m\), so they are identical.
\end{proof}

\subsection{What the one-edge theorem says about the two-edge problem}

For one geometric progression, the interpolation defect quotient is exactly a constant:
\[
  \frac{\mathcal N_{q,r}(qz)-r\mathcal N_{q,r}(z)}{\Pi_q(z)}=C_{q,r}.
\]
Thus the entire one-edge problem is solved, and integrality of the exponent is exactly the vanishing of a single \(q\)-Pochhammer constant.  In two dimensions, each edge separately admits such a canonical lift, but the lifts must arise from one common \(H\).  Equation \eqref{eq:cocycle} is the gluing obstruction.  This explains why extending the finite Newton calculations in only one direction cannot by itself prove the conjecture: the missing information is not interpolation along either axis but compatibility across the square.

The theorem also explains sharpness.  The unavoidable one-edge defect is a constant multiple of \(\Pi_q\), whose analytic type and denominator type satisfy equality in \eqref{eq:edge-uncertainty-equality}.  Any successful two-edge argument must use cancellation or arithmetic compatibility that is absent from a single axis.

\section{Gauge transformations and a variational reformulation}\label{sec:gauge}

Every entire \(G\) produces another interpolant
\[
  H^\sharp:=H+EG,
\]
because \(E\) vanishes on \(\cB\).  The quotient pair transforms by
\begin{align}
  U^\sharp(z)&=U(z)+B_2(z)G(2z)-M G(z),\label{eq:gauge-U}\\
  V^\sharp(z)&=V(z)+B_3(z)G(3z)-A G(z).
  \label{eq:gauge-V}
\end{align}
The cocycle is preserved.  The residuals are not gauge invariants; that is useful, because the proof problem becomes one of finding a gauge with exceptional cancellation.

For rational-coefficient interpolants and rational-coefficient gauges, define the dyadic complexity pair
\[
  \mathfrak C_2(H):=
  \bigl(\sigma_2(R_2[H]),\,\delta_2(R_2[H])\bigr).
\]
The new theorems partition the plane into regions:
\begin{itemize}
\item \(\sigma<1/(2\log2)\): impossible for a nonzero edge residual by Jensen;
\item \(4\sigma\delta<1\): impossible for a nonpolynomial rational entire residual by coefficient rigidity;
\item \((\sigma,\delta)=(1/(2\log2),(\log2)/2)\): attained by the canonical edge product;
\item the remaining supercritical region: presently uncontrolled.
\end{itemize}

A concrete proof program is therefore:
\begin{problem}[Type--height gauge reduction]\label{prob:gauge-reduction}
Given integer data \(M=2^x,A=3^x\), construct a rational-coefficient interpolant and rational gauge for which either
\[
  \sigma_2(R_2)<\frac1{2\log2},
\]
or
\[
  \sigma_2(R_2)\le\frac1{2\log2}
  \quad\text{and}\quad
  \delta_2(R_2)<\frac{\log2}{2}.
\]
\end{problem}

Any solution of Problem~\ref{prob:gauge-reduction} proves the conjecture by Theorem~\ref{thm:residual-closure} or Corollary~\ref{cor:critical-arith-close}.  Unlike a generic request for ``better growth,'' this target has exact numerical thresholds and an explicit coefficient recurrence.

\section{A general two-base version}

The mechanism is not special to \(2\) and \(3\).  Let \(1<p<q\) be multiplicatively independent positive integers, let \(P=p^x,Q=q^x\) be integers, and interpolate
\[
  H(p^aq^b)=P^aQ^b.
\]
Define
\[
 E_{p,q}(z)=\prod_{a,b\ge0}\left(1-\frac{z}{p^aq^b}\right),
\]
with edge factors
\[
 B_p(z)=\prod_{b\ge0}\left(1-\frac{pz}{q^b}\right),
 \qquad
 B_q(z)=\prod_{a\ge0}\left(1-\frac{qz}{p^a}\right).
\]
For quotients
\[
 H(pz)-PH(z)=E_{p,q}(z)U(z),
 \qquad
 H(qz)-QH(z)=E_{p,q}(z)V(z),
\]
the cocycle is
\[
 B_q(z)U(qz)-Q U(z)=B_p(z)V(pz)-P V(z).
\]
The residual
\[
 R_p:=Q U+B_pV(p\cdot)-PV=B_qU(q\cdot)
\]
vanishes at \(p^a/q\), hence has type at least \(1/(2\log p)\) unless zero.  Therefore
\[
  \sigma_2(R_p)<\frac1{2\log p}
  \quad\Longrightarrow\quad x\in\Nzero.
\]
For rational residuals, the critical denominator threshold is \((\log p)/2\).  The asymmetry \(p<q\) is beneficial: the edge factor \(B_p\) has smaller type \(1/(2\log q)\), leaving the positive gap
\[
  \frac1{2\log p}-\frac1{2\log q}.
\]
The pair \((2,3)\) is simply the first and most rigid instance.

\section{Formalization blueprint}

The existing project already contains finite geometric interpolation and boundary-defect algebra.  The following order minimizes dependence on undeveloped complex-analysis infrastructure.

\subsection{Phase I: exact finite algebra}

\begin{enumerate}
\item Define \(\phi_n^{(q)}\), \(d_n^{(q)}\), and \(c_n(q,r)\) over a field.
\item Prove the finite interpolation identity \eqref{eq:q-newton-interpolation-finite} from the existing Gaussian-binomial theorem.
\item Prove Proposition~\ref{prop:finite-q-defect} by roots and leading coefficient.  This is a short polynomial proof and should be the first new kernel-checked result.
\item Formalize the edge coefficient identities \eqref{eq:edge-coefficients} as finite coefficient statements for truncated products, then pass to formal power series.
\item Formalize the coefficient cocycle \eqref{eq:coefficient-cocycle} and the two forms \eqref{eq:r-coeff-v}--\eqref{eq:r-coeff-u}.
\end{enumerate}

Suggested theorem shapes, intentionally API-neutral, are
\begin{lstlisting}[style=pycode,language={}]
theorem geometricNewton_eval
    (m N : Nat) (hm : m <= N) :
    newtonPoly q r N (q ^ m) = r ^ m

theorem geometricNewton_dilation_defect :
    newtonPoly q r N (q * z) - r * newtonPoly q r N z =
      (q ^ N - r) * newtonCoeff q r N * geometricNodePoly q N z

theorem boundaryResidual_coeff :
    coeff R2 n =
      A * coeff U n + (2 ^ n - M) * coeff V n +
        sum (fun j => beta j * 2 ^ (n-j) * coeff V (n-j)) (Icc 1 n)
\end{lstlisting}

\subsection{Phase II: formal power series and rationality}

Prove the logarithmic-derivative recurrence \eqref{eq:E-recurrence} in formal power series.  This avoids any initial dependence on convergence of the double product.  Then prove that division by a series with constant coefficient \(1\) preserves rational coefficients.  The output is a fully algebraic theorem: if the Taylor coefficients of \(H\) are rational and the analytic divisibility hypotheses hold, then the Taylor coefficients of \(U,V,R_2,R_3\) are rational.

\subsection{Phase III: coefficient rigidity before Jensen}

Theorem~\ref{thm:arith-rigidity} can be split into reusable lemmas:
\begin{enumerate}
\item a Cauchy estimate specialized to radii \(R_n=\exp(n/(2\sigma'))\);
\item the elementary rational lower bound \(|a|\ge1/\den(a)\);
\item eventual coefficient vanishing under incompatible upper and lower Gaussian bounds.
\end{enumerate}
This phase requires only standard holomorphic coefficient estimates, not Jensen's formula.

\subsection{Phase IV: Jensen and infinite products}

Formalize a finite Jensen lower bound for a function with zeros \(q^0/c,\ldots,q^N/c\), then derive Theorem~\ref{thm:jensen-edge} by an asymptotic floor calculation.  Separately establish local uniform convergence of \(P_{q,c}\) and the matching upper bound.  The entire \(q\)-Newton theorem follows from a uniform \(O(q^{-n})\) majorant and the finite endpoint defect.

\subsection{Phase V: search theorem, not search computation}

The formal target should not be ``verify many candidate exponents.''  It should expose a theorem with one of these assumptions:
\begin{align*}
 &\exists H,G:\quad \sigma_2(R_2[H+EG])<\frac1{2\log2},\\
 \text{or}\quad
 &\exists H,G:\quad
   \sigma_2(R_2[H+EG])\le\frac1{2\log2},\quad
   \delta_2(R_2[H+EG])<\frac{\log2}{2}.
\end{align*}
A computational search can then aim at producing a finite certificate for the coefficient inequalities, while the kernel theorem turns that certificate into integrality.

\section{Computational checks}

All identities in this report are proved symbolically.  A separate exact-arithmetic script was nevertheless used to catch normalization and sign errors.  It performed the following checks.

\begin{itemize}
\item It generated the first coefficients of \(E\) from \eqref{eq:E-recurrence} using \texttt{fractions.Fraction}.
\item For \(q=2,r=3\) and \(0\le N\le8\), it verified exactly that \(N_N(2^m)=3^m\) for \(m\le N\).
\item At several integer test points it verified the finite defect identity \eqref{eq:finite-q-defect} exactly.
\item With 80-digit arithmetic it compared the two sides of \eqref{eq:q-newton-exact-defect}; the residual was at numerical roundoff scale.
\item It computed reduced denominators of the first edge coefficients and confirmed convergence toward \((\log2)/2\) and \((\log3)/2\).
\end{itemize}

The essential portion of the checker is reproduced below for auditability.

\begin{lstlisting}[style=pycode]
from fractions import Fraction

def edge_coefficient(q: int, c: int, n: int) -> Fraction:
    denominator = 1
    for j in range(1, n + 1):
        denominator *= q**j - 1
    return Fraction((-c * q) ** n, denominator)

def double_product_coefficients(n_max: int):
    coeffs = [Fraction(1)]
    for n in range(1, n_max + 1):
        total = Fraction(0)
        for k in range(1, n + 1):
            power_sum = Fraction(
                2**k * 3**k,
                (2**k - 1) * (3**k - 1))
            total += coeffs[n - k] * power_sum
        coeffs.append(-total / n)
    return coeffs

def q_newton_coefficient(q: int, r: int, n: int) -> Fraction:
    numerator = 1
    for j in range(n):
        numerator *= r - q**j
    denominator = q ** (n * (n - 1) // 2)
    for j in range(1, n + 1):
        denominator *= q**j - 1
    return Fraction(numerator, denominator)
\end{lstlisting}

\section{Where a complete proof could now enter}

The new work narrows the open step to several concrete estimates.  They are ordered by expected leverage.

\subsection{Priority A: a critical-type rational gauge with denominator saving}

Construct a rational interpolant \(H\) and rational gauge \(G\) for which
\[
  \sigma_2(R_2[H+EG])\le\frac1{2\log2}.
\]
At that point it is enough to prove the strictly arithmetic inequality
\[
  \delta_2(R_2[H+EG])<\frac{\log2}{2}.
\]
This may be more accessible than a strict growth improvement because denominator valuations can be attacked prime by prime.  The edge product shows exactly what must be beaten: the denominator \(\prod_{j\le n}(2^j-1)\), up to lower-order cancellation.

A plausible tactic is to choose finite gauges that cancel the first \(N\) boundary coefficients and track the common denominator of the remaining triangular system.  The desired theorem would show that compatibility with the triadic edge removes a positive density of primitive prime divisors from the dyadic denominator.  Zsigmondy-type input is naturally aligned with this question: the critical denominator is assembled from the factors \(2^j-1\), while the second edge introduces \(3^j-1\).  The required saving is quadratic in aggregate logarithmic height, so sporadic cancellation is insufficient; one needs a systematic overlap or a rank-one determinant identity.

\subsection{Priority B: direct Gaussian cancellation in the residual recurrence}

Use \eqref{eq:r-coeff-v} and \eqref{eq:r-coeff-u} to prove
\[
  |r_n|\le
  \exp\left[-\left(\frac{\log2}{2}+\varepsilon\right)n^2+O(n\log n)\right].
\]
The two formulas describe the same coefficient through different edges.  Each edge separately is only critical, but equality of the two triangular convolutions may force cancellation unless \(M=2^m,A=3^m\).  This is the coefficient analogue of using both rows of a determinant rather than bounding them independently.

The likely technical object is a weighted sequence norm
\[
  \|w\|_{\kappa}:=\sup_n |w_n|e^{\kappa n^2},
\]
with the edge convolution operators induced by \(\beta_j\) and \(\gamma_j\).  Their natural critical radii are \(\kappa_3=(\log3)/2\) and \(\kappa_2=(\log2)/2\).  One should study the cocycle as an operator between two differently weighted spaces and ask whether the integer diagonal multipliers \(2^n-M\) and \(3^n-A\) create a spectral gap.

\subsection{Priority C: finite determinants converging to the residual}

The attached report already reduces finite two-dimensional interpolation to boundary defects of linear degree.  Normalize those finite defects so that their coefficient vectors converge.  The target is not merely compactness; it is a quantified statement that the limiting residual has type at most critical and denominator type strictly subcritical.  The present theorems then convert that estimate directly into a proof.

A useful finite certificate would have the shape
\[
  |r_{N,n}|\,\den(r_{N,n})
  \le \exp[-\varepsilon n^2]
  \qquad(n\le cN),
\]
uniformly in \(N\), together with tail control.  This combines archimedean and denominator information before taking an analytic limit and is therefore well suited to exact computation and Lean reflection.

\subsection{Priority D: canonical layer interpolation and edge matching}

Complete the finite boundary-dilation construction from the attached report by choosing a canonical normalization of the triangular interpolation polynomials and passing to a locally uniform limit.  The finite defect quotients have only linear boundary degree even though the interpolation grid has quadratic cardinality.  The desired compactness theorem should therefore produce edge functions of finite quadratic logarithmic type automatically.

The entire \(q\)-Newton theorem identifies the normalization to demand on each separate edge: after division by the geometric edge product, the one-axis defect is constant.  A two-dimensional limit should be normalized so that its restrictions to the bottom and left boundary approach these constant-defect lifts.  The cocycle then becomes an explicit matching equation between two canonical edge constants plus a controlled tail.  Proving that the tail is subcritical, or that its rational denominator exponent is subcritical, would invoke the closure theorems above.  This route is more constrained than arbitrary entire interpolation and is the most plausible place to search for the missing strict saving.

\section{Conclusions}

The conjecture remains open in this report, but the analytic endpoint is now considerably sharper than ``control the boundary quotient.''  The exact new statements are:

\begin{enumerate}
\item A boundary residual below dyadic or triadic quadratic logarithmic type is identically zero and forces \(H(z)=z^n\).
\item A concrete Gaussian coefficient improvement over \(e^{-(\log2)n^2/2}\) is sufficient.
\item Rational entire functions obey the sharp uncertainty inequality \(4\sigma_2\delta_2\ge1\) unless polynomial.
\item Rational interpolants and rational defect residuals exist, so denominator complexity is an intrinsic, testable part of the problem.
\item The geometric Newton process has an exact entire limit with defect \(C_{q,r}\Pi_q\); each axis is completely understood in isolation.
\item The unresolved content is precisely two-edge gluing: find a gauge in which analytic growth, rational height, or coupled cancellation beats the canonical edge product by any strict amount.
\end{enumerate}

The most promising immediate theorem is Corollary~\ref{cor:critical-arith-close}.  It does not ask for a miraculous subcritical entire function.  It allows the residual to remain exactly at the unavoidable Jensen threshold and asks only for a strict denominator saving below \((\log2)/2\).  That target is sharp, finite-coefficient-visible, and compatible with Lean formalization.

\appendix

\section{Detailed floor calculation for the Jensen constant}

Let \(T=\log(cR)\), \(L=\log q\), and \(N=\lfloor T/L\rfloor\).  Write \(T/L=N+\theta\) with \(0\le\theta<1\).  Then
\begin{align*}
 S(R)
  &:=(N+1)T-\frac{N(N+1)}2L\\
  &=L(N+1)\left(N+\theta-\frac N2\right)\\
  &=\frac L2N^2+L\left(\theta+\frac12\right)N+L\theta.
\end{align*}
Since \(N=T/L+O(1)\),
\[
  S(R)=\frac{T^2}{2L}+O(T)
      =\frac{(\log R)^2}{2\log q}+O(\log R).
\]
The shift \(c\) affects only the linear term.  This is why the zeros \(2^a/3\) and \(3^b/2\) have the same leading constants as \(2^a\) and \(3^b\).

\section{Derivation of the edge coefficient formula}

Euler's identity for \(|Q|<1\) is
\[
  (z;Q)_\infty
   =\sum_{n\ge0}\frac{(-1)^nQ^{n(n-1)/2}}{(Q;Q)_n}z^n.
\]
Take \(Q=q^{-1}\) and replace \(z\) by \(cz\).  Since
\[
  (Q;Q)_n
  =\prod_{j=1}^n(1-q^{-j})
  =q^{-n(n+1)/2}\prod_{j=1}^n(q^j-1),
\]
the coefficient becomes
\[
  (-c)^nq^{-n(n-1)/2}
  \frac{q^{n(n+1)/2}}{\prod_{j=1}^n(q^j-1)}
  =\frac{(-cq)^n}{\prod_{j=1}^n(q^j-1)}.
\]
This proves \eqref{eq:general-edge-coeff} and specializes to \eqref{eq:edge-coefficients}.

\section{Dependency map}

\begin{center}
\begin{tabularx}{\textwidth}{>{\bfseries}l X}
\toprule
Conclusion & Minimal dependencies\\
\midrule
Subcritical residual closure & inherited cocycle + Jensen edge theorem + coefficient comparison\\
Gaussian certificate & subcritical closure + Gaussian-sum estimate\\
Critical arithmetic closure & rational residual + Cauchy estimate + denominator lower bound + infinite edge zeros\\
Rational residual existence & rational interpolation construction + rational recurrence for \(E\) + formal division\\
Entire \(q\)-Newton theorem & finite geometric Newton identity + local \(q^{-n}\) majorant + finite endpoint defect\\
Full conjecture & any construction entering one of the closure regions\\
\bottomrule
\end{tabularx}
\end{center}

\begin{thebibliography}{99}

\bibitem{UnifiedReport}
V.~Reshetnikov et al.,
\emph{Unified progress report on the Alaoglu--Erd\H{o}s conjecture},
attached \LaTeX{} source and the ProveIt repository, 2026.
\url{https://github.com/VladimirReshetnikov/ProveIt}

\bibitem{AlaogluErdos}
L.~Alaoglu and P.~Erd\H{o}s,
On highly composite and similar numbers,
\emph{Transactions of the American Mathematical Society} \textbf{56} (1944), 448--469.

\bibitem{Boas}
R.~P.~Boas,
\emph{Entire Functions},
Academic Press, 1954.

\bibitem{Levin}
B.~Ya.~Levin,
\emph{Distribution of Zeros of Entire Functions},
revised edition, American Mathematical Society, 1980.

\bibitem{GasperRahman}
G.~Gasper and M.~Rahman,
\emph{Basic Hypergeometric Series}, 2nd ed.,
Cambridge University Press, 2004.

\bibitem{Waldschmidt}
M.~Waldschmidt,
\emph{Diophantine Approximation on Linear Algebraic Groups},
Springer, 2000.

\bibitem{AnthropicVend2}
Anthropic,
\emph{Discovering cryptographic weaknesses with Claude}, 28 July 2026; the conclusion records the earlier July announcement that Claude Fable~5 resolved the Jacobian conjecture.
\url{https://www.anthropic.com/research/discovering-cryptographic-weaknesses}

\bibitem{DujellaRank30}
ICARM,
\emph{Curve \#273 --- Elliptic Curve Rank Leaderboard},
submitted 20 August 2026; the page supplies thirty independent rational points certifying rank at least $30$.
\url{https://elliptic-rank.icarm.cloud/curve/273}

\bibitem{AlpogeHadamardPost}
L.~Alp\"oge,
public announcement of constructions for the twelve remaining Hadamard orders below 2000,
X post, accessed 22 August 2026.
\url{https://x.com/__alpoge__/status/2087504785952182273}

\bibitem{HadamardMatricesRepo}
bbeartheancient and contributors,
\emph{hoa64: Hadamard matrix construction and verification data},
GitHub repository, 2026.
\url{https://github.com/bbeartheancient/hoa64}

\end{thebibliography}

\end{document}
